\section{Data and pipeline}\label{sec:data}

\begin{figure*}[tb]
\centering
\resizebox{\textwidth}{!}{% The pipeline, sources to paper. Hand-laid-out (a diagram, not a measurement);
% every count on it is a macro from numbers.tex / numbers2.tex.
\begin{tikzpicture}[
  x=1cm, y=1cm,
  box/.style={draw=ink, rounded corners=1.5pt, fill=white, inner sep=2.5pt,
              align=center, font=\scriptsize, minimum height=7mm, text width=2.55cm},
  src/.style={box, fill=paleblue},
  store/.style={box, fill=palegold},
  derive/.style={box, fill=palegrey},
  outbox/.style={box, fill=white, draw=ink, thick},
  arr/.style={-{Latex[length=1.5mm]}, draw=ink!75, line width=0.5pt},
  lab/.style={font=\tiny, fill=white, inner sep=1pt, text=ink!80},
]
% ---- row 1: sources
\node[src] (igu)  at (1.5, 4.3) {IGU fleet list\\ \tiny annual PDF};
\node[src] (ais)  at (4.7, 4.3) {AISstream\\ \tiny 3 WebSockets per worker, 150 MMSI slots, 1\,h rotation};
\node[src] (vf)   at (7.9, 4.3) {VesselFinder API\\ \tiny masterdata; credit-budgeted rescue};
\node[src] (noaa) at (11.1, 4.3) {NOAA Marine Cadastre\\ \tiny US Class-A AIS, 2016--25};
\node[src] (gfw)  at (14.3, 4.3) {Global Fishing Watch\\ \tiny port visits and voyages, 2016--26};
% ---- row 2: stores
\node[store] (reg) at (1.5, 2.7) {\texttt{vessel\_registry}\\ \tiny \nFleet{} LNG carriers and FSRUs};
\node[store] (fix) at (6.3, 2.7) {\texttt{ais\_fixes}\\ \tiny \nFixes{} fixes (hypertable)};
\node[derive] (sm) at (9.5, 2.7) {state machine\\ \tiny per vessel, \nZonePolygons{} polygons at \nTerminals{} terminals};
\node[store] (pe) at (12.7, 2.7) {\texttt{port\_events}\\ \tiny \nEvents{} events, regime-tagged};
\node[src] (mkt) at (15.9, 2.7) {EIA, FRED, AGSI+,\\ ERA5, TTF quote\\ \tiny prices, storage, weather};
% ---- row 3: derivations
\node[derive] (wl) at (3.9, 1.1) {\texttt{priority\_watchlist}\\ \tiny 5 tiers, re-scored hourly};
\node[derive] (pair) at (9.5, 1.1) {legs, visits, queues\\ \tiny the three pairings};
\node[store] (sd) at (12.7, 1.1) {\texttt{signal\_daily}\\ \tiny \nSignalKeys{} keys, two bases, \nSignalRows{} rows};
\node[store] (mp) at (15.9, 1.1) {\texttt{model\_panel}\\ \tiny \nPanelFeatures{} features, publication-lagged};
% ---- row 4: analysis and paper
\node[outbox] (parta) at (9.5, -0.5) {Part A replays\\ \tiny A1, A2, A4, A5 nowcasts};
\node[outbox] (partb) at (12.7, -0.5) {Part B replays\\ \tiny M0--M2 ladder, FWL scan, H2--H4};
\node[outbox] (paper) at (15.9, -0.5) {tables, macros, figures\\ \tiny generated, never typed};
% ---- edges
\draw[arr] (igu) -- (reg);
\draw[arr] (vf) -- node[lab, pos=0.62] {IMO$\to$MMSI} (reg);
\draw[arr] (ais) -- (fix);
\draw[arr] (vf) -- node[lab, pos=0.55] {rescue fixes} (fix);
\draw[arr] (noaa) -- node[lab, pos=0.5] {backfill} (fix);
\draw[arr] (fix) -- (sm);
\draw[arr] (sm) -- (pe);
\draw[arr] (gfw) -- node[lab, pos=0.5] {events written directly} (pe);
\draw[arr] (reg.south) |- (wl.west);
\draw[arr] (fix.south west) -- node[lab, pos=0.5, right=1pt] {last fix, ETA, destination} (wl.north east);
\draw[arr] (wl.north) -- node[lab, pos=0.22, left=1pt] {subscription list} (ais.south);
\draw[arr] (pe) -- (pair);
\draw[arr] (pair) -- (sd);
\draw[arr] (sd) -- (mp);
\draw[arr] (mkt) -- (mp);
\draw[arr] (pair) -- (parta);
\draw[arr] (mp) -- (partb);
\draw[arr] (parta) -- (partb);
\draw[arr] (partb) -- (paper);
\end{tikzpicture}
}
\caption{The pipeline from sources to paper. Blue boxes are external sources, gold
boxes are stored tables, grey boxes are derivations, and the outlined boxes are the
replays that generate every table and figure. The watchlist closes a loop: the
ranking computed from stored fixes decides which vessels the live feed subscribes
to next.}
\label{fig:pipeline}
\end{figure*}

\subsection{Sources}

Vessel positions come from three feeds of different fidelity, and no signal is
aggregated across the boundary between them. For the United States we use the NOAA
Marine Cadastre archive of terrestrial Class-A AIS \citep{noaa_ais}, which is
exhaustive along the US coast from 2016 and ends on \noaaEnds. For Europe and the
open ocean we use Global Fishing Watch port-visit and voyage events \citep{gfw}.
These supply arrivals and departures but carry no coordinates and no anchorage
events; the feed ends on \gfwEnds. A live subscription to an AIS aggregator covered a
tier-ranked fleet from April~2026 (Section~\ref{sec:live}) and is used only for the
vintage self-validation described in Section~\ref{sec:panel}.

The fleet in scope is every LNG carrier and floating storage and regasification
unit (FSRU) in the IGU fleet list, \nFleet{} hulls, resolved to an MMSI through a
vessel-database lookup. The signals concern the transatlantic subset of their
movements: laden departures from US export terminals, arrivals and discharges at
European import terminals, and the ballast return. Cargoes bound for Asia or Latin
America enter the panel only as an unresolved-destination band
(Section~\ref{sec:partb}).

Prices and controls are public. \HH{} spot and Lower-48 storage are from the EIA API
\citep{eia_api}. The TTF front-month contract is from a keyless quote feed and was
cross-checked against the World Bank's monthly European gas series
\citep{worldbank_cmo}; the regression slope is 0.999 with $\Rsq=0.9999$ over 106
months, so the distortion introduced by the continuous roll is negligible. EUR/USD
and Brent are from FRED \citep{fred}. EU storage is from GIE AGSI+ \citep{agsi}.
Heating and cooling degree days are computed from ERA5 reanalysis temperatures at
weighted demand centres \citep{openmeteo}.

\subsection{Reconstruction}\label{sec:reconstruction}

Table~\ref{tab:inventory} gives the scale of the reconstruction and
Figure~\ref{fig:pipeline} its shape. Fixes are resolved against \nZonePolygons{}
polygons at \nTerminals{} terminals in seven zones (US Gulf, US Atlantic, north-west
Europe, Baltic, Iberia, western Mediterranean and eastern Mediterranean). The
polygons were drawn by hand in QGIS from satellite imagery and port charts, one
berth, one anchorage and one approach envelope per terminal, with sub-zones where a
terminal has several berths. The approach envelope is deliberately drawn to contain
the anchorage, the channel and the berth, so that a single visit stays open while a
vessel transits between them; without it the four-to-six hour channel transit at
Sabine Pass or Rotterdam would register as a departure and a fresh arrival.
Figure~\ref{fig:zones} shows the three polygons at Sabine Pass with a sample of the
events emitted inside them.

\begin{table*}[tb]
\centering\small
\caption{Reconstruction inventory. Port events are labelled by the feed that produced
them. The two backfill feeds end months before the live feed, which determines the
sample choice in Section~\ref{sec:partb}.}
\label{tab:inventory}
\begin{tabular}{lrl}
\toprule
layer & rows & span \\
\midrule
AIS position fixes & 54,494,686 & 2016-01-01 to 2026-08-10 \\
port events (state machine) & 87,874 & \\
  \quad noaa & 44,363 & to 2025-12-31 \\
  \quad gfw & 40,273 & to 2026-02-21 \\
  \quad mmsi\_filter & 1,974 & to 2026-08-10 \\
  \quad bbox & 1,264 & to 2026-05-30 \\
in-scope fleet (LNG carriers + FSRUs) & 828 & \\
laden US export departures & 9,855 & \\
signal panel rows (\texttt{signal\_daily}) & 1,032,982 & 34 keys, two bases \\
model panel rows (\texttt{model\_panel}) & 77,500 & 20 features \\
\bottomrule
\end{tabular}

\end{table*}

\begin{figure}[t]
\centering
\includegraphics[width=\linewidth]{fig7_zones}
\caption{The hand-drawn polygons at \zoneFigureTerminal{} and a random sample of
the \nZoneFigureEvents{} coordinate-bearing events the state machine emitted there.
The approach envelope (light) contains the anchorage (hatched) and the berth (gold,
at the head of the channel). Zone entries scatter along the envelope boundary,
anchored events cluster in the anchorage boxes, and moored and departed events sit
on the berth.}
\label{fig:zones}
\end{figure}

\begin{figure*}[tb]
\centering
\resizebox{0.78\textwidth}{!}{% The per-vessel state machine (pipeline/state_machine.py). Thresholds are the
% code's defaults: 30 min dwell at SOG < 1 kn for anchored/moored, 15 min at
% SOG > 1 kn outside the berth for departed. Events fire back-dated to the first
% qualifying fix, not to the fix that confirmed the dwell.
\begin{tikzpicture}[
  x=1cm, y=1cm,
  st/.style={draw=ink, rounded corners=2pt, fill=palegrey, font=\scriptsize\bfseries,
             inner sep=3pt, minimum width=1.9cm, minimum height=6mm},
  ev/.style={font=\tiny, align=center, fill=white, inner sep=1pt},
  arr/.style={-{Latex[length=1.5mm]}, draw=ink, line width=0.55pt},
]
\node[st] (T) at (0, 0) {TRANSIT};
\node[st] (E) at (4.2, 0) {IN\_ENVELOPE};
\node[st] (A) at (8.8, 1.9) {ANCHORED};
\node[st] (M) at (8.8, -1.9) {MOORED};
\node[st] (D) at (13.2, -1.9) {DEPARTED};
\draw[arr] (T) -- node[ev, below, yshift=-2pt] {\texttt{zone\_entry}\\first fix inside\\the approach envelope} (E);
\draw[arr] (E) -- node[ev, below right, pos=0.55, xshift=-2pt, yshift=-1pt] {\texttt{anchored}\\$\ge$30\,min at SOG$<$1\,kn\\in the anchorage} (A);
\draw[arr] (E) -- node[ev, below left, pos=0.5, xshift=2pt, yshift=-1pt] {\texttt{moored}\\$\ge$30\,min at SOG$<$1\,kn\\in a berth} (M);
\draw[arr] (A) -- node[ev, right] {\texttt{moored}} (M);
\draw[arr] (A) to[bend right=38] node[ev, above, yshift=3pt] {\texttt{anchorage\_exit}} (E);
\draw[arr] (M) -- node[ev, above, yshift=2pt] {\texttt{departed}\\SOG$>$1\,kn outside\\the berth for 15\,min} (D);
\draw[arr] (D) to[bend left=30] node[ev, below, yshift=-3pt] {\texttt{zone\_exit}: first fix outside the envelope} (T);
\node[ev, text=ink!70] at (4.2, 3.0) {\texttt{anchorage\_entry} / \texttt{anchorage\_exit}: raw polygon crossings, no dwell filter};
\end{tikzpicture}
}
\caption{The per-vessel state machine. Each fix is first resolved to a polygon
(berth wins over anchorage, which wins over approach; an open visit is sticky to its
terminal), then the state is advanced. Dwell-confirmed events are back-dated to the
first qualifying fix. Transitions require a fix outside the relevant polygon, so a
gap in coverage never emits an event on its own.}
\label{fig:statemachine}
\end{figure*}

A per-vessel state machine (Figure~\ref{fig:statemachine}) walks each vessel's
fixes in time order and emits zone entry, anchorage, mooring and departure events.
A vessel is \term{anchored} after at least 30 minutes inside an anchorage polygon at
speed over ground below 1\,kn, \term{moored} after the same dwell inside a berth
polygon, and \term{departed} after 15 minutes outside the berth at more than 1\,kn.
Anchorage entries and exits are raw polygon crossings with no dwell filter, so
that queue times can be measured without the half-hour back-dating of the
dwell-confirmed marker. A vessel whose first observed fix is already inside a
polygon receives synthetic entry and mooring events flagged as a cold start.

Laden state is inferred from the draught the vessel broadcasts relative to its
design draught,
\begin{equation}
  \text{laden}_e = \mathbf{1}\!\left[\, d(t_e) \ge 0.85\, d_{\mathrm{design}} \,\right],
  \label{eq:laden}
\end{equation}
where $d(t_e)$ is the most recent reported draught at or before the event time
for inbound events and, for departures, the first draught reported in a window
from one hour before to six hours after the event, because a master typically
rebroadcasts the lighter draught only after undocking. When no usable draught
exists the state is set from the terminal's flow direction: a vessel that moored at
an export terminal leaves laden, and one that moored at an import terminal leaves
in ballast.

Three pairings of those events underlie everything that follows and are drawn in
Section~\ref{sec:panel}. A \term{leg} pairs a departure with the vessel's next zone
entry and represents a voyage. A \term{visit} pairs a mooring with the next departure
and represents berth occupancy. A \term{queue} pairs the first anchorage entry with
the next mooring at the same terminal and represents the wait to berth.

\subsection{Live acquisition}\label{sec:live}

The live feed is the part of the project that took most of the engineering, and its
limits are part of the result, so it is described here even though the decade
panel is what the tests run on. The AIS aggregator caps a subscription at 50 vessel
identifiers per connection. The system therefore runs two ingester workers on
separate egress addresses, each holding three WebSocket connections: two connections
form a persistent block of 100 slots and the third is a scan rotation of 50 slots
that is re-drawn every hour, for 300 slots in total against a fleet of \nFleet.

Which vessels hold a slot is decided by a ranking recomputed hourly from stored
fixes and static data (Figure~\ref{fig:watchlist}). Tier~1 is any vessel seen
inside a terminal polygon in the last three days. Tier~2 is a vessel with a parsed
ETA within 48 hours or a declared destination that resolves to a known terminal.
Tier~3 is a vessel inside the wider zone bounding box, ordered by whether it is
closing on a terminal. Tiers~4 and~5 are recently and not recently seen. Two
classes are pinned into the persistent block regardless of tier: a vessel on an
open leg inside its expected arrival window in either direction, and a vessel
currently in port, so that a long berth queue does not decay out of coverage.
FSRUs are demoted to the lowest tier because a deployed FSRU sits moored for
months and its own fixes never drive a signal. At the time of writing
\nWatchlistSlotted{} of \nWatchlistVessels{} ranked vessels hold a slot:
\nSlotsPinned{} pinned, \nSlotsPersistent{} persistent and \nSlotsScan{} in
rotation.

\begin{figure}[t]
\centering
\includegraphics[width=\linewidth]{fig9_watchlist}
\caption{Subscription slots by tier at the time of writing. The horizontal axis is
symmetric-log so that the \nSlotsNone{} unslotted vessels in tier~5 do not hide the
allocation in tiers~1 to~4.}
\label{fig:watchlist}
\end{figure}

Terrestrial AIS does not see the open ocean, and a vessel that leaves a slot goes
dark until it is re-ranked. Two backstops exist. First, a vessel that goes silent
near a terminal while a leg-defining event is at risk (a laden departure, an import
arrival, a berthing) is looked up on a commercial vessel-position API and the
returned position is injected as an ordinary fix, so the state machine
re-acquires it for free. The lookups cost credits from a reserve that expires on a
fixed date, so the daily budget is derived each run as the straight line from the
current balance to zero at expiry; leg-defining classes are exempt from the line and
lower classes spend only the surplus above it (Figure~\ref{fig:rescue}). Between
the first snapshot and \vfCreditsAsOf{} the reserve fell from \vfCreditsStart{} to
\vfCreditsNow{} credits against an expiry of \vfExpiry, over \nRescueLookups{}
lookups of which \nRescued{} returned a usable position; \nRescueSkipped{} candidate
lookups were declined by the budget, which is the measured unmet demand. Second, a
daily discovery pass resolves newly delivered hulls from the fleet list, because
under a closed subscription list no vessel can otherwise enter the system.

\begin{figure*}[tb]
\centering
\includegraphics[width=\textwidth]{fig10_rescue}
\caption{The credit-budgeted position backstop. Left: outcomes of every lookup
attempt by rescue class, on a symmetric-log axis; green is a usable position, grey
is a lookup that returned nothing, and light grey is a candidate the budget declined.
Right: the account balance against the glide line from the first snapshot to zero at
expiry, which is the policy the daily cap is derived from.}
\label{fig:rescue}
\end{figure*}

Two limits follow directly. A vessel that was never ranked highly enough to hold
a slot leaves no live trace at all, and a vessel that drops out of coverage mid-voyage
is only recovered if it re-enters a slot or is bought back. Both are absent from the
decade archives, which is why the tests of Sections~\ref{sec:parta} and~\ref{sec:partb}
run on the NOAA and GFW backfills and the live feed serves only to validate the
point-in-time construction.

\subsection{Ground truth}\label{sec:groundtruth}

The US side of the reconstruction can be checked against a published total. The EIA
reports monthly US LNG exports $X_m$ in MMcf of gas. Dividing by a nominal cargo of
$V_c = 174{,}000$\,m$^3$ of LNG, which at a volumetric expansion of 600 is about
3.69\,Bcf of gas, gives an implied cargo count, and the ratio of observed laden
departures $D_y$ to that count is the capture rate for year $y$:
\begin{equation}
  c_y \;=\; \frac{D_y}{\sum_{m\in y} X_m \,/\, (600\,V_c)} .
  \label{eq:capture}
\end{equation}
Restricted to the NOAA feed, $c_y$ is \captureTwentyTwenty{} in 2020 and rises to
\captureTwentyThree{}, \captureTwentyFour{} and \captureTwentyFive{} in 2023, 2024
and 2025 (Table~\ref{tab:capture}, Figure~\ref{fig:capture}). The gradient is a
property of the archive, not of the vessels: the NOAA archive is assembled from
terrestrial receivers whose density along the Gulf coast grew over the decade, and
a cargo that loaded in 2020 out of range of a receiver left no fix. The GFW feed has
the opposite gradient, strong in the years the NOAA archive is thin and fading after
2022, because it is fed by satellite receptions whose coverage of the fleet
declined. The two feeds are not fused in this paper. The project's diagnostic tool
counts every feed and therefore double-counts departures seen by both, and a deduped
union calibrated to the EIA total is listed as future work
(Section~\ref{sec:future}).

\begin{table}[t]
\centering\small
\caption{NOAA-only capture of US LNG export cargoes relative to the EIA-implied count,
Equation~\eqref{eq:capture}.}
\label{tab:capture}
\begin{tabular}{lrrr}
\toprule
year & captured & EIA-implied & capture \\
\midrule
  2016 & 48 & 51 & 95\% \\
  2017 & 159 & 192 & 83\% \\
  2018 & 216 & 294 & 74\% \\
  2019 & 290 & 494 & 59\% \\
  2020 & 300 & 648 & 46\% \\
  2021 & 715 & 966 & 74\% \\
  2022 & 918 & 1048 & 88\% \\
  2023 & 1,099 & 1178 & 93\% \\
  2024 & 1,216 & 1185 & 103\% \\
  2025 & 1,570 & 1494 & 105\% \\
\bottomrule
\end{tabular}

\end{table}

\begin{figure}[t]
\centering
\includegraphics[width=\linewidth]{fig19_capture}
\caption{Observed laden departures against the EIA-implied cargo count by year. The
capture rate (gold, right axis) is the coverage of the NOAA archive, which is why
early-year levels of every US signal are biased low.}
\label{fig:capture}
\end{figure}

The consequence for modelling is stated here because it matters everywhere below.
Early-year levels of any US signal are biased low by up to half. The models
therefore never treat an early-year level as a physical quantity. The physical
nowcasts (Section~\ref{sec:parta}) forecast within-window changes and score against
nulls computed from the same observed series, so a uniform under-count deflates the
model and the null identically. The spread scan (Section~\ref{sec:partb})
residualises each signal on the controls and on its own past before it is used, so a
slowly varying level bias is absorbed by the persistence term. What neither design
can absorb is a level-dependent non-linearity in the early years, and that is listed
as a limitation.

Figure~\ref{fig:spread} shows the target and the headline signal on one timeline.
The two structural events that any model must survive are visible: the 2020 and
2021 inversions, in which \HH{} briefly exceeded \TTF{}, and the 2022 European
crisis.

\begin{figure*}[tb]
\centering
\includegraphics[width=0.92\textwidth]{fig1_spread_and_stock}
\caption{The \HH{}--\TTF{} spread (top) and EU-bound gas at sea (bottom), daily, on
the coverage-consistent sample used for modelling. The extremes coincide with Winter
Storm Uri (February~2021) and the August~2022 TTF peak. The collapse of the at-sea
stock in the final weeks is not a physical event: it is the point at which the two
backfill feeds end and only the thin live feed remains (Section~\ref{sec:seam}),
which is why the primary sample is truncated there.}
\label{fig:spread}
\end{figure*}
